%!TEX root = ../../thesis.tex

\section{Constraints Imposed by the Learner}\label{sec:mdp-constraints}

% TODO: This is the load-bearing section --- the design decisions in
% \cref{sec:input-actions,sec:corpus-actions} are consequences of it, so state the
% constraints crisply and cite the implementation:
%
%   C1  Single flat categorical action. The action fed to the dynamics network is a
%       one-hot of width Action::size(), and MCTS branches over 0..action_count
%       using the index directly. Any composite action must be packed into one
%       integer.
%   C2  No action masking. legal_actions is asserted to be the identity 0..action_size
%       (efficientzero learner/mcts.rs) --- subset or phase-dependent legal sets are
%       not wired end to end. Every encodable action must be legal in every state.
%   C3  Fixed-width observation. Observation::to_vec() has a constant length, so the
%       corpus, the coverage map and any per-slot features must be encoded at a fixed
%       size.
%   C4  Deterministic transitions. MCTS plans over learned latent dynamics and
%       therefore wants a deterministic (state, action) -> next_state. Stochastic
%       mutation operators break the planner's model, which is why
%       \cref{sec:corpus-determinism} replaces randomness with an exhaustive sweep.
%   C5  Branching factor. The policy head is one softmax of width action_size and
%       MCTS enumerates every child, so the action space must stay enumerable
%       (order 10^3), not merely representable.
%
% Note explicitly that C4 and C5 are the price of choosing a planning algorithm over
% a model-free one such as the DQN of \cite{bottinger_deep_2018} --- they are forced,
% not stylistic.
